\noindent Let $p_1=2<p_2=3<p_3=5<\cdots$ be the primes, and fix an integer $k\ge 2$. Write
\[
P_k := \prod_{i=1}^k p_i
\]
(the $k$-th primorial). The question asks whether there exists $N(k)$ such that for all integers $n\ge N(k)$ there exists an integer
\[
 m\in [n,n+P_k)
\]
with ``$>k$ many prime factors.''

\noindent \emph{Ambiguity note (distinct vs. counted with multiplicity).}
If ``prime factors'' were counted \emph{with multiplicity} ($\Omega(m)$), then for $n>P_k$ the claim would be immediate: the interval $[n,n+P_k)$ contains a multiple $M$ of $P_k$, and $M\ge 2P_k$, so $\Omega(M)\ge \Omega(P_k)+1 = k+1$. This would make the case $k=2$ with interval length $6$ trivial, contradicting the remark in the problem text that the case $k=2$ is unknown. Therefore the intended meaning must be:

\begin{quote}
$\omega(m)>k$, i.e. $m$ has more than $k$ \emph{distinct} prime divisors.
\end{quote}

\noindent Under this interpretation, the question becomes:

\noindent \textbf{Conjecture 891(k).} For each $k\ge 2$ there exists $N(k)$ such that for all $n\ge N(k)$, the interval $[n,n+P_k)$ contains an integer $m$ with $\omega(m)\ge k+1$.

\noindent The text also notes: (a) Schinzel deduced from a theorem of P\'olya about gaps between $k$-smooth numbers that Conjecture 891(k) holds with $P_k$ replaced by $p_1\cdots p_{k-1}p_{k+1}$; (b) Weisenberg observed that Dickson's conjecture would make Conjecture 891(k) false if $P_k$ is replaced by $P_k-1$.

\noindent 2)
